倒排文件（IVF）索引将数据库划分为多个列表，通常按向量最近的质心进行分配。查询搜索其中 $n_{\rm probe}$ 个列表。若 $\mathcal L(q)$ 表示选中列表的编号集合，则可参与评分的向量为
\[
 U(q)=\bigcup_{\ell\in\mathcal L(q)}\mathrm{list}_\ell,\qquad n_c=|U(q)|.
\]
位于 $U(q)$ 之外的真实近邻，在向量评分开始之前就已被排除。

乘积量化（product quantisation，PQ）将向量坐标划为 $M$ 个子空间，为每个子空间学习码本，从而压缩向量 \cite{pq}。在子空间 $m$ 中，编码 $z_m$ 指向代表向量 $\mathcal C_m[z_m]$；这些代表向量共同重建 $\widehat x$。非对称距离计算（asymmetric distance computation，ADC）保留未量化的查询，计算
\begin{equation}
 \|q-\widehat x\|^2=\sum_{m=0}^{M-1}\|q_m-\mathcal C_m[z_m]\|^2.
 \label{eq:2-adc}
\end{equation}
其中 $q_m$ 是查询的子向量。每条查询先计算到各码字的距离并制成表，随后扫描数据库编码时使用查表。PQ 分离了不同子空间，但任意学习得到的码本未必能进一步分离同一子空间\emph{内部}的坐标。JHQ 的主码本具有这一额外结构。

精排对初步评分选出的候选短名单投入更多计算 \cite{johnson_gpu,faiss_refine}。令 $\alpha$ 为精排倍率，$c=\max\{k,\lceil\alpha k\rceil\}$ 为请求的短名单大小。在 $|U(q)|\geq k$ 时，以 $C_\alpha(q)$ 表示至多 $c$ 个保留候选，以 $\mathcal R_\alpha(q)$ 表示最终返回的 $k$ 个向量，则
\begin{equation}
 \mathcal R_\alpha(q)\subseteq C_\alpha(q)\subseteq U(q)\subseteq\mathcal X.
 \label{eq:2-cascade}
\end{equation}
若 $\mathcal T_k$ 为该查询的 GT 集合，由包含关系可得
\[
 |\mathcal R_\alpha\cap\mathcal T_k|
 \leq |C_\alpha\cap\mathcal T_k|
 \leq |U\cap\mathcal T_k|.
\]
增加 probe 数会扩大路由候选池；增加 $\alpha$ 会扩大请求的短名单。要区分路由、粗筛或重构导致的召回损失，需要在保持其他阶段固定的条件下逐项比较。下一节定义 JHQ 的两级评分。
